% commentaire.tex — f171r, vv. 721–764 (conversion initiale des travaux antérieurs) ; discussions textuelles déplacées dans critique.tex (voir repartition.md)
\comm{721–5}{ἀθέοις ἔργοις γὰρ ἀνταμει- / ψόμεσθά σ᾽ ὥσπερ εἰκὸς ἀντὶ τῶνδε / τάχα δὲ \{σε\} μεταβαλοῦσ᾽ ἐπὶ κακὸν ἑτερότροπον / ἐπέχει \{τις\} τύχη.}{Le grand Jean Taillardat avait donné un exposé sur son interprétation des vers 721-725 du manuscrit de Ravenne. Il nous offre une explication limpide de ce passage fort débattu ; je transcris pieusement ci-dessous le résumé de son intervention, tout en reprenant sa traduction\nc{À retrouver dans la \italique{Revue des Études Grecques}, tome 101, fascicule 482-484 (juillet-décembre 1988). Actes de l'association, p. XVIII-XIX.} : \\
« Car toi, avec tes impiétés, nous te châtierons de ta conduite, comme de juste. Et puis la Fortune, changeant brusquement d'humeur, aura vite fait de poursuivre un méchant homme qui a pris le chemin du crime ».\\
Cette traduction appelle les justifications suivantes : \\
— Ἀθέοις ἔργοις : selon Jean Taillardat, il est impossible d'interpréter ce datif comme le complément de moyen ἀνταμειψόμεσθα et de traduire « \latin{\italique{Impiis factis} enim remunerabimus te sicut par est pro eis} » (Andrea Divo) ou « \latin{\italique{Nefariis} enim \italique{factis} vicissim, ut par est, haec tua ulciscemur} » (Brunck) ou « Nous y répondrons, comme de raison, \italique{par des actes sacrilèges} » (Willems)\nc{Ou « \anglais{Since your crimes are godless, it is fitting that our punishment of you be \italique{godless} as well} » (Austin \& Olson.)}. En effet, on ne pouvait pas considérer comme sacrilège « \italique{la coutume} (qui) était d'apporter du feu près des suppliants réfugiés sur l'autel » afin de les contraindre à le lâcher (Σ Eur. \italique{Androm}. 256 : ἔθος ἦν πῦρ προσφέρειν τοῖς εἰς βωμὸν καταφεύγουσιν ; cf. Eur. \italique{Herc. fur}. 243, 285). Ἀθέοις ἔργοις est donc : \ \ \ 1) \ soit un datif-instrumental d'accompagnement lié à σ(ε), cf. Kühner-Gerth, I, p. 435; ἀθέοις ἔργοις a, dans ce cas, la même fonction syntaxique que ξὺν ἀδίκοις ἔργοις (v. 716) lié à σοι (v. 715) ; \ \ \ 2) \ soit un datif-instrumental de cause lié à ἀνταμειψόμεσθα. Jean Taillardat penche pour la première solution.\\
— Ἀντὶ τῶνδε (v. 723) $=$ ἀντὶ τῶνδε τῶν ἀθέων ἔργων « nous te payerons en retour de tes impiétés».\\
— Μεταβαλοῦσ(α), v. 724, est employé absolument, ce qui n'est pas rare, cf. L.S.J. \italique{s. v.} μεταβάλλω, A. III. 2.\\
On reconnaîtra d'autre part la locution ἐπέχειν ἐπί τινα (Hérodote 9.59) à peu près synonyme de διώκειν τινά (Id. \italique{ibid}.), car τάχα ... ἐπὶ κακὸν ἑτερότροπον ἐπέχει (présent de prophétie) τύχη équivaut à τάχα ... ἐπὶ κακὸν ἄνδρ᾽ ὃς τὴν ἑτέραν ὁδὸν ἐτράπετ' ἐπέχει τύχη.\\
Pour κακός $=$ κακός τις ἀνήρ, cf. Théognis 189, 190, 577, 972, 1112 (de même κακὴ $=$ κακή τις γυνή, id. 185). Le composé ἑτερότροπος « qui a pris la mauvaise route » (avec ἕτερος euphémique pour κακός, cf. L.S.J. III. 2) est une allusion au mythe bien connu des deux chemins, celui du bien et celui du mal (Hésiode \italique{Op}. 287-292 ; Prodicos in Xénophon \italique{Mém}. 2. 1. 21-27). \\
Quoi qu'on en ait dit, le sens de ce passage est très clair : le chœur fait entendre au Parent d'Euripide qu'à défaut d'une punition donnée par les femmes, il recevra un châtiment infligé par le Ciel.\\
Pour le texte, voir le commentaire critique.}
\comm{724–5}{τάχα δὲ \{σε\} μεταβαλοῦσ᾽ ἐπὶ κακὸν ἑτερότροπον ἐπέχει \{τις\} τύχη.}{Il suffit que le chœur dise `τάχα δὲ' ... et le rythme même du vers s'accélère : les trois dochmiaques des vers 724-725 contiennent force résolutions.}
\comm{726–7}{ἀλλὰ τάσδε μὲν λαβεῖν χρῆν σ᾽ ἐκφέρειν τε τῶν ξύλων, / καὶ καταίθειν τὸν πανοῦργον πυρπολεῖν θ᾽ ὅσον τάχος.}{C'était, en effet, l'habitude de faire périr par le feu ceux qui avaient cherché refuge sur un autel. Cf. Euripide, \italique{Androm}., v. 256 ; \italique{Herc. Fur}., v. 243 et 285 et les interprètes. (Note de Willems.)\\
Cette \italique{exhortation} -- κατακελευσμός selon la dénomination consacrée de Zieliński\nc{Voir \italique{The Verse of Greek Comedy} (1912), p. 316, de John W. White. Le lecteur pourra utilement consulter ce livre de référence, entre autres s'il s'intéresse à la structure des comédies d'Aristophane.} -- de la Coryphée pour appeler à la reprise de l'action est selon l'usage un distique de tétramètres trochaïques. (Et, comme au vers 531-532, débute par `ἀλλά'.)}
\comm{726}{Τάσδε}{renvoie à des compagnes de Mica et non pas à des torches, selon l'hypothèse invraisemblable de Fritzsche -- qui a malheureusement fait florès. Logique pratique élémentaire: avant d'aller recueillir du bois, il faut d'abord se \italique{débarrasser} de ce qui encombre -- en l'occurrence : les torches -- et non s'en embarrasser. Ce n'est qu'une fois le bois entassé que Mica pourra effectivement se saisir d'une torche et bouter le feu au bûcher.}
\comm{730}{τὸ Κρητικὸν.}{Hésychius (κ 4087) : \citComm{Κρητικόν· ἱματίδιον λεπτὸν καὶ βραχύ.\\
\francais{Crétique : petit habit fin et court.}}}
\comm{733}{ἀσκὸς ἐγένεθ᾽ ἡ κόρη.}{Il était probablement interdit d'apporter du vin lors des Thesmophories, fêtes consacrées à Déméter et sa fille. (Cf. par exemple l'adage latin \latin{Cereri sacrificant}\nc{Recensé par Érasme, \italique{Adages} n° 3071. Voir aussi les adages n° 1895, 1896.}, \italique{sacrifier à Cérès}, c'est-à-dire faire abstinence de vin.) Cette interdiction théorique devait souffrir quelques écarts pratiques, si l'on en croit le compte-rendu de Mnésiloque (vers 626-633).}
\comm{734}{Περσικὰς ἔχων.}{On comprend d'après d'autres passages d'Aristophane que les \italique{persiques} désignent des chaussures légères et recouvrant le pied\nc{La puce des socratiques expériences des \italique{Nuées}, v. 151, se retrouve chaussée de `persiques' après que Socrate lui ait plongé les pattes dans la cire – voir aussi Hésychius (π 2005) : \grec{περσικά· εὐτελῆ ὑποδήματα}.}, peut-être des \italique{chaussons}\nc{Les Athéniennes pouvaient les garder au lit, en prenant garde toutefois de maintenir les jambes en l'air, cf. \italique{Lysistrata} v. 229.} ; en tout cas, ces \italique{persiques} sont portées par les femmes\nc{\italique{Assemblée des femmes} v. 319 et Hésychius (μ 929) : \grec{μεσοπερσικαί· ὑποδήματα γυναικεῖα}.}, ou comme ici les enfants.\\
Les «~persiques~» dont parle Mnésiloque sont visibles sur le fameux cratère de Würzbourg\nc{Reproduit en annexe. Voir la note aux vers 754-757.}. Il faut savoir que les outres à l'époque étaient faites de la peau d'un animal évidé\nc{Pour plus d'informations, aller consulter l'indispensable Daremberg \& Saglio, \italique{Dictionnaire des antiquités grecques et romaines} à l'entrée \latin{Uter}.}, bouc ou chèvre la plupart du temps. Ces «~persiques~» sont en réalité les pattes arrières de l'animal outré, fermées par un lacet.}
\comm{735}{ὦ θερμόταται γυναῖκες.}{θερμός évoque l'\italique{impétuosité} féminine : \anglais{hot-headed, hasty}, l'\italique{audace} des femmes\nc{Cf. le \LSJ , Taillardat §§313, 351.}.
Le mot peut aussi désigner la \italique{fièvre} ; il n'y a cependant pas d'équivalent grec à notre \italique{fièvre de Bercy} : à ma connaissance, le grec n'a jamais rapproché la fièvre de l'ivrognerie.
Quand donc le scholiaste à notre vers écrit `ὦ μέθυσοι', il glose `ὦ ποτίσταται' et non `ὦ θερμόταται'.
Suidas (θ 251) s'y est trompé qui recopie la glose en la faisant précéder du lemme θερμόταται :\citComm{Θερμόταται: μέθυσοι. ὦ θερμόταται γυναῖκες, μέγα καπήλοις καλόν, ἡμῖν δὲ κακόν [...]}Le cas n'est pas isolé : Suidas avait déjà commis un précédent du même genre : au vers 178 (\italique{confer} le commentaire \italique{ad locum}), lisant κοινῇ à la place de καινῇ, il n'hésita pas à utiliser la glose afférente, μεγάλῃ καὶ θαυμαστῇ, comme nouvelles acceptions de κοινῇ (en y ajoutant le véritable sens de κοινῇ : ἀνθρωπίνῃ).
Ce sont là des méthodes qui montrent que la \italique{Souda} est un lexique à manipuler avec précaution.}
\comm{738}{κακὸν δὲ καὶ τοῖς σκευαρίοις καὶ τῇ κρόκῃ.}{Mnésiloque reproche à ces soûlaudes d'Athéniennes de contribuer à la ruine des ménages par leur inconduite biberonnesque.
Montrons-le !
\grec{τό σκευάριον}, diminutif de σκεῦος, désigne littéralement un \italique{petit récipient} (\italique{menue vaisselle}) ou \italique{petit meuble}.
Employé au pluriel, il peut désigner plus particulièrement le meuble où l'on entasse son argent. \italique{Confer} le passage du \italique{Ploutos} ci-dessous (vers 806-811) qui nous renseigne sur l'usage du mobilier grec :\citComm{Ἡ μὲν σιπύη μεστή ᾽στι λευκῶν ἀλφίτων,\\
οἱ δ᾽ ἀμφορῆς οἴνου μέλανος ἀνθοσμίου.\\
Ἅπαντα δ᾽ ἡμῖν ἀργυρίου καὶ χρυσίου\\
\textbf{τὰ σκευάρια} πλήρη ᾽στίν, ὥστε θαυμάσαι.\\
Τὸ φρέαρ δ᾽ ἐλαίου μεστόν· αἱ δὲ λήκυθοι\\
μύρου γέμουσι, τὸ δ᾽ ὑπερῷον ἰσχάδων.\\
\francais{Notre huche au pain regorge d'une blanche farine, /
nos amphores d'un vin noir qui embaume. /
Tous nos \italique{meubles}} (τὰ σκευάρια) \francais{sont pleins /
d'argent et d'or, que c'est merveille. /
La citerne regorge d'huile, les lécythes /
sont emplis d'essences, le grenier de figues sèches}}
(C'est bien de \italique{pièces} d'or et d'argent dont il est question puisque les plats et autres assiettes d'airain ou d'argent sont évoqués aux vers suivants.)
Il est tout naturel que ces meubles (coffrets ?) servant de porte-monnaie aient pris par métonymie le sens de \italique{biens}, \italique{économies}.
C'est d'ailleurs si naturel que c'est effectivement le cas : aux vers 838-839, le Juste, contant ses misères passées, emploie τὰ σκευάρια en ce sens métonymique :
\citComm{Δίκαιος — κομιδῇ μὲν οὖν· / αὐχμὸς γὰρ ὢν \textbf{τῶν σκευαρίων} μ᾽ ἀπώλεσεν.\\
\francais{Le Juste — Précisément. Une sécheresse s'étant mise dans \italique{mon mobilier} ($=$ mes économies) fut ce qui me perdit.}}Cette acception de τὰ σκευάρια est encore plus manifeste dans un passage au début des \italique{Grenouilles}. Dionysos y cherche un moyen d'expédier en Enfer ses \italique{bagages} ; son esclave Xanthias regimbant, il hèle un mort en route pour l'Hadès (vers 170-172) :
\citComm{Καὶ γάρ τιν᾽ ἐκφέρουσι τουτονὶ νεκρόν.\\
Οὗτος, σὲ λέγω μέντοι, σὲ τὸν τεθνηκότα.\\
Ἄνθρωπε, βούλει \textbf{σκευάρι᾽} εἰς Ἅιδου φέρειν;\\
\francais{Car voilà justement un mort qu'on emporte.\\
Holà ! c'est à toi que je parle, toi le mort.\\
Homme, veux-tu porter à l'Hadès quelques \italique{petits colis} ?}}Il y a un évident détournement dans la plus pure veine aristophanesque de l'universel adage : on n'emporte (φέρειν) pas son \italique{argent} (σκευάρι᾽) dans l'autre monde (εἰς Ἅιδου) ; les σκευάρια de l'adage sont à prendre au sens figuré, \italique{argent} mais se retrouvent tout matériellement incarnés en \italique{bagages} dans la version illustrée d'Aristophane\nc{Même procédé dans les \italique{Nuées} : quand Strepsiade désabusé songe au \italique{Tourbillon} (Δῖνος) de Socrate, il ne voit plus en cette nouvelle divinité qu'une vile \italique{poterie} (δῖνος). Si l'on veut savoir ce que le véritable Socrate entendait par ce \italique{Tourbillon}, il suffit d'interroger un autre grand esprit socratique : Pierre Boutang : «~La lignée [de race et d'esprit], c'est merveilleux !.. mais un peu trop linéaire pour mon goût. Donc il me faut quelque chose de plus circulaire : des lignes qui se répètent, qui reviennent sur elles-mêmes, qui se retournent ; des tourbillons ! [ ... ] Un grand tourbillon créateur !~». Quoi qu'on en ait dit, le Socrate croqué par Aristophane est loin d'être si caricatural et malgré ses défauts qu'Aristophane grossit à la loupe, il est encore à mille coudées au-dessus de Cléon, par exemple.}. Voilà pour τοῖς σκευαρίοις, passons à τῇ κρόκῃ.\\
\grec{ἡ κρόκη} est au propre \italique{le fil de trame}. 
Il possède, comme nous l'allons voir, le même sens imagé de \italique{richesses} que τὰ σκευάρια.
Le grec, comme le français, a prodigalement tiré maintes images et maints proverbes de l'activité filandière et tisserande. Observons l'un d'eux : σπαθᾶν τὴν κρόκην.
Au vers 55 des \italique{Nuées}, Strepsiade énumère les oppositions entre lui, le rustre économe, et sa femme, la raffinée dilapidatrice. Il pousse la comparaison jusque dans l'alcôve où sa femme souffre de besoins tyranniques ; il lui reproche de \italique{tirer trop sur la corde} -- le grec dit \italique{frapper sur le fil} :\citComm{ὦ γύναι, λίαν σπαθᾷς.}σπαθᾶν [τὴν κρόκην] désigne l'action de battre chaque fil de trame (κρόκη) au moyen de la σπάθη, ou \italique{battant}, pour rendre le tissu compact\nc{Avec le commentaire du scholiaste repris dans la \italique{Souda} (ε 3176) : καὶ γὰρ ἐπὶ τῶν ὑφαινομένων ἱματίων λέγομεν σπαθᾶν τὸ ἄγαν κρούειν τὴν κρόκην. Taillardat (§440) développe l'exégèse de l'image. Voir aussi le lexique de Photius (ε 2014) : Ἐσπαθημένος· πεπυκνωμένος· ἀπὸ τῶν κατακρουουσῶν σπαθῶν τὴν κρόκην.} ; ce faisant, l'on utilise davantage de laine, d'où le fait que l'expression signifie \italique{dissiper l'argent}, \italique{gaspiller}.
Or l'on trouve les expressions σπαθᾶν τὰ πατρῷα\nc{Diphil. 43, 27.}, τὰ χρήματα\nc{Plut., \italique{Per.} 14. Il s'agit ici des \italique{deniers publics} : τῶν δὲ περὶ τὸν Θουκυδίδην ῥητόρων καταβοώντων τοῦ Περικλέους ὡς σπαθῶντος τὰ χρήματα καὶ τὰς προσόδους ἀπολλύντος.}, τάλαντα\nc{Lucien, \italique{Cat.} 20.}. 
Le sème commun à ces différents mots est celui de \italique{richesse} (matérielle ou pécuniaire) ; puisqu'ils ont été substitués à τὴν κρόκην, c'est donc la preuve -- taillardesque\nc{Lire et relire l'excellentissime article du regretté Jean, « Images et matrices métaphoriques » (1977), in \italique{Bulletin de l'Association Guillaume Budé} 4, pp. 344-354.} -- qu'ils partagent avec κρόκη ce sème.
Par exemple, en français l'on rencontre les expressions \italique{ourdir un complot}, \italique{une intrigue}, \italique{une conspiration}, \italique{une vengeance}. Pris dans leur contexte, le sème commun aux différents compléments d'ourdir est : \italique{machination secrète}. Or l'expression tisserandière propre est «~ourdir une trame~», qui consiste à entrecroiser les fils de trame avec les fils ourdis, c'est-à-dire disposés sur l'ourdissoir. Les termes de complot, intrigue, \&c. ne sont substituables à la trame que parce qu'ils ont en commun avec elle le sème de \italique{machination secrète} (sens figuré de trame).\\
Finissons avec un amusant proverbe sur ce fléau qu'est l'alcool, surtout pour le sexe filandier\nc{\italique{Corpus paroemiographorum Graecorum}, Appendices, III, 3.} :\citComm{ἡ γραῦς τὰς ἀγαθίδας ἐκπίνουσα\\
La vieille boit\nc{Sur l'image de «~manger~» son bien, voir Taillardat §439.} ses pelotes.}ἡ ἀγαθίς est une \italique{pelote de fil}. On constate que ce mot a aussi pris le sens de \italique{bien}, \italique{richesse} (on parlerait d'un \italique{bas de laine} en français) ; par permutation métaphorique avec κρόκη ? peut-être aussi par paronomase : ἀγαθῶν ἀγαθίδες est une expression désignant une grande quantité de biens.
Notre tradition populaire n'est d'ailleurs pas en reste : voici ce que l'on chantait à l'époque où l'on chantait encore\nc{Doit-on dire de cette époque qu'«~elle est révolue, qu'elle a fait son temps ; qu'on n'la reverra plus... et qu'c'est bien attristant~» ? Toujours est-il que j'emprunte ces chansons à une tradition qui a depuis longtemps défunté, mais qui fut néanmoins -- et heureusement ! -- bien embaumée par notre grand Patrice Coirault, aux cotes 110-05 et 110-06 de son \italique{Répertoire des chansons françaises de tradition orale}.} :\citComm{\francais{Parlons un peu des filandières.\\
Celles qui filaient l'étoupe\\
Ont mangé toute la soupe.\\
Celles qui filaient le lin\\
Ont lampé tout le vin.}}Et encore, à propos de filles qui se sont enivrées :\citComm{\francais{Elles ne l'ont pas fait en filant du lin,\\
Mais elles l'ont bien fait en buvant du vin.}}}
\comm{741}{δέκα μῆνας.}{Les mois lunaires grecs étant de 29 ou 30 jours, la femme portait bien environ dix mois l'enfant. Aulu-Gelle dans ses \italique{Nuits attiques} traite pendant tout un chapitre (III, 16) de la question du temps de gestation. N'en déplaise au scholiaste, la durée de 10 mois est bien attestée, notamment par Ménandre ; Aulu-Gelle cite aussi Varron expliquant que les noms des Parques \italique{Nona} et \italique{Decima} viennent des nombres de mois fixés pour l'accouchement (\latin{item Nona et Decima a partus tempestivi tempore}).}
\comm{743}{τρικότυλον ἢ πῶς;}{C'est-à-dire du vin supérieur, à une obole les trois cotyles. Hésychius (τ 1388) : \citComm{\grec{τρικότυλος οἶνος· οὗ τρεῖς κοτύλαι ὀβολοῦ πωλοῦνται},}ce qui est cher, en effet, car cela met le métrète (144 cotyles, environ 40 litres) à $7fr68$\nc{$7fr68$ de feu notre franc germinal~; c'est-à-dire qu'on pouvait échanger à tout moment contre $7,68 \times 4,5 = 34,56$ grammes d'argent pur~; soit effectivement 48 oboles (de $0,72$ grammes dans le système éginète)... mais cela, c'était avant l'invention des monnaies de singe : chez nous, bien-sûr, et déjà en $-406$ à Athènes (\italique{confer} les \italique{Grenouilles}, vers 718-733).}, le double à peu près du prix ordinaire. Démosthène (1045.5) désigne de la même manière un vin plus cher encore, \grec{ὁ οἶνος δωδεκάδραχμος}, le vin à douze drachmes, c'est-à-dire à $11fr52$ le métrète\nc{72 oboles soit $51,84$ grammes d'argent.}. À rapprocher de τετρώβολον, désignant dans la \italique{Paix}, 254, le miel attique. [...] — Ἤ πῶς; est mis pour ἤ πῶς ὤνιος; Cf. \italique{Acharn}., 758, \italique{Caval}., 480. (Note de Willems.)\\
Mnésiloque s'était déjà adressé à Mica en lui reprochant d'avoir acheté son enfant au vers 565.\\
Pour le texte, voir le commentaire critique.}
\comm{746}{πόσ᾽ ἔτη δὲ γέγονε; τρεῖς Χοᾶς ἢ τέτταρας;}{S'il faut en croire le scholiaste, l'on pouvait compter l'âge à partir du nombre de Grandes Dionysies écoulées (\grec{τοῦτο ἔθος ἦν Ἀττικοῖς, λέγειν τὰ ἔτη κατὰ τὸν ὑπερπίπτοντα ἀριθμὸν ἀπὸ τῶν Διονυσίων}) ; usage qui correspond chez nous à celui de compter les \italique{printemps}. Mnésiloque, filant la métaphore de l'outre et l'enfant, préfère se référer malicieusement aux \italique{Choés}, journée la plus arrosée des Anthestéries\nc{L'on trouvera d'abondants et utiles renseignements sur cette fête à l'article \latin{Dionysia} du \italique{Dictionnaire des antiquités grecques et romaines} de Daremberg \& Saglio. Cette fête des Anthestéries correspond à peu près à nos Saint-Vincent, Percée du Vin Jaune ou autres fêtes du beaujolais nouveau.} ; d'autant qu'un \grec{χοῦς} est aussi une unité de volume\nc{Le \grec{χοῦς} – ou \italique{conge} romain — correspond à 12 cotyles, soit environ 3 litres ¼.}. Cela permet insidieusement à Mnésiloque de préciser que l'outre n'est pas si petite : elle est d'une capacité d'une dizaine de litres\nc{D'une capacité \italique{théorique} de trois ou quatre conges : il est évident que Mica ne s'est pas encombrée du fardeau que serait cette outre pleine (Mnésiloque aurait d'ailleurs eu le plus grand mal à la lui chaparder). Une telle capacité peut étonner mais il faut bien voir que ces outres étaient faites de la peau complète d'un animal (souvent un bouc ou une chèvre) : ainsi, Strepsiade (\italique{Nuées}, v. 1238), trouvant que son créancier ferait une belle outre, l'estime de six conges. Il y avait des outres d'une capacité impressionnante pour jouer au jeu de l'\grec{ἀσκωλιασμός}, consistant à sauter -- \italique{clocher}, comme l'étymologie l'indique -- et tenir en équilibre sur une outre graissée. Voir l'article \latin{Uter} du \italique{Dictionnaire des antiquités grecques et romaines} de Daremberg \& Saglio. Il y a de toute façon une exagération comique : l'outre représentée sur le cratère de Würzbourg n'a certainement pas cette capacité et Mnésiloque est dans l'outrance verbale.}.}
\comm{747}{σχεδὸν τοσοῦτον χὤσον ἐκ Διονυσίων.}{L'on traduit généralement ce vers par quelque chose d'équivalent peu ou prou à : « C'est cela, plus le temps écoulé depuis les Dionysies ». Outre qu'inventer un anniversaire à l'outre est une billevesée indigente relevant davantage de l'humour des philologues que de celui d'Aristophane, il présente une impossibilité : essayez donc de conserver pendant trois ou quatre \italique{printemps} une outre de vin dans les conditions de l'Athènes de l'an $-2300$ avant le Réfrigérateur !
Et puisque Mica ne peut pas parler de l'âge de son outre, elle parle donc de sa \italique{capacité}. En l'occurrence, elle essaye de se défendre du chef d'accusation de biberonne porté par Mnésiloque~: sa chère petite outre, «~peut-être est-elle de cette taille [c'est-à-dire d'une capacité d'une dizaine de litres] mais elle est telle [c'est-à-dire aussi peu remplie] depuis les Dionysies~». Autrement dit, Mica au jour des Thesmophories prétend ne pas avoir touché au contenu de l'outre depuis les dernières Dionysies, ce qui représente une abstinence de sept mois ! (Bluff !)}
\comm{748}{μὰ τὸν Ἀπόλλω τουτονί.}{Contre les commentateurs qui se contentent d'une banale formule de ce genre pour inventer sur scène la présence d'une statue d'Apollon omniprésente,
Willems avait déjà répondu par une note au vers 83 des \italique{Nuées}. Comme ce genre d'idées saugrenues se retrouve encore dans les dernières éditions les plus autorisées, j'ai cru faire œuvre utile en reproduisant ladite note en annexe.}
\comm{754–7}{Scène d'égorgement de l'outre nourrissonne.}{Cette scène est illustrée sur le cratère de Würzbourg\nc{Cratère apulien d'environ 370 av. J.-C., Martin von Wagner Museum der Universität Würzburg, Antiken-Sammlung, H5697. L'on pourra se reporter pour plus de détails à l'article d'Eric Csapo, \italique{A Note on the Würzburg Bell-Crater H5697} («~Telephus Travestitus~») in \italique{Phoenix}, Vol. 40, No. 4 (Winter, 1986), p. 379-392.} (reproduit en annexe). Amusante mise en abîme : le `σφαγεῖον' que tient Mica sur cette représentation n'est autre qu'un cratère de forme identique au cratère de Würzbourg lui-même -- c'est-à-dire un vase propre à recevoir du vin pour l'y mélanger et non une coupe servant à recevoir le sang des victimes (σφαγεῖον).}
\comm{755–7}{τό γ᾽ αἷμα τοῦ τέκνου τοὐμοῦ λάβω --- ὕπεχ᾽ αὐτό, χαριοῦμαι γὰρ ἕν γε τοῦτό σοι --- κακῶς ἀπόλοι᾽, ὡς φθονερὸς εἶ καὶ δυσμενής.}{La mise en scène de ces vers n'a pas bien été comprise\nc{Les commentateurs avaient imaginé jusqu'ici que : \ \ soit Mnésiloque déversait le contenu de l'outre sur l'autel ; \ \ soit il buvait d'un trait le vin. Cette dernière option est notamment l'opinion du scholiaste : τοῦτο ἅμα εἰπὼν ἐκχέει τὸν οἶνον ὁ κηδεστὴς ἐν τῷ βωμῷ. Je passe sur les autres explications alambiquées du même acabit. Austin \& Olson, conformément à leur pente, ne se prononcent pas et avancent les deux hypothèses. Il faudrait cesser de croire que la pusillanimité est une preuve d'esprit critique. «~La recherche de la vérité n'empêche pas de prendre parti~», comme disait l'autre imbécile.}.
Mica a demandé à \italique{récupérer} le sang de son enfant (τό γ᾽ αἷμα τοῦ τέκνου τοὐμοῦ λάβω) et puisque Mnésiloque lui concède cette unique faveur (χαριοῦμαι γὰρ ἕν γε τοῦτό σοι), il faut l'en croire.
Seulement, Mnésiloque prend au mot Mica. Le \italique{mot}, en l'occurrence, c'est λάβω : \italique{prendre}, \italique{récupérer}, dans la bouche de Mica. Mnésiloque entend \italique{recevoir}. C'est entendu : elle n'a qu'à \italique{présenter} (ὕπεχ᾽ αὐτό) le σφαγεῖον. Il lui jette alors à la figure le contenu de l'outre. C'est précisément cet instant, où Mica tend le σφαγεῖον, qui est saisi sur le cratère de Würzbourg. On voit qu'il suffit à Mnésiloque de balancer le bras pour asséner l'outre à fragmentation liquide à la tête de sa mère.
C'est ce comportement dépassant, outrepassant la mesure -- μηδὲν ἄγαν, Mnésiloque, μηδὲν ἄγαν. Ne te l'a-t-on pas assez répété ? --, qui explique l'emportement de Mica : κακῶς ἀπόλοι᾽, ὡς φθονερὸς εἶ καὶ δυσμενής. Devant Mica maculée du «~sang~» de sa fille, ajoutez la confusion induite par les nombreux quiproquos\nc{Le grec dirait qu'on surajoute \italique{du feu au feu}, πῦρ ἐπὶ πῦρ.}, et vous comprendrez aisément les quelques précisions que demande Critylla déboussolée à son arrivée au vers 759 : `τί τῆς ἱερείας γίγνεται;' puis `ταλαντάτη Μίκα τίς ἐξεκόρησέ σε;'.
(L'exégèse de ce passage est d'ailleurs la preuve \italique{in texto} que Mica présente des signes \italique{visibles} d'agression -- agression qui n'est donc pas uniquement psychologique. Ces signes, comme nous l'allons voir à la note s'y rapportant, sont les taches de vin qui imprègnent son vêtement.)}
\comm{758}{τουτὶ τὸ δέρμα τῆς ἱερείας γίγνεται.}{τὸ δέρμα désigne l'outre vide – qui rappelons-le était faite d'une peau d'animal évidé. L'usage était que le prêtre reçoive certaines parties de la victime (\grec{τά ἱερώσυνα}).}
\comm{759}{}{Pour le texte, voir le commentaire critique.\\
Faut-il prendre au premier degré la déclaration de Mnésiloque et en inférer que c'est Critylla qui jouait le rôle de prêtresse ?}
\comm{760–1}{τίς ἐξεκόρησέ σε; / τίς τὴν ἀγαπητὴν παῖδά σου ᾿ξῃρήσατο;}{Faute d'avoir compris que Mnésiloque venait d'asperger Mica du «~sang~» de sa propre fille, ce passage n'a pas non plus été entendu.
Même le grand Willems -- une fois n'est pas coutume -- fut pris la main dans le sac à ordures : trompé par un jeu de mots en latin\nc{\latin{Quis te devirginavit} ?} qu'il a cru pouvoir rapprocher de ce distique, il proposa comme interprétation une paillardise qui n'a pas lieu d'être, et qui d'ailleurs n'y est pas.\\
Figurons-nous la scène : l'outre a été égorgée et déversée sur Mica ; arrive Critylla à qui Mnésiloque jette l'outre vide. Critylla avec l'«~enfant~» éviscéré ; Mica maculée de sang face à elle.\citComm{Critylla --- τίς ἐξεκόρησέ σε;}Le verbe ἐκκορέω, \italique{balayer}, est fréquemment utilisé comme synonyme de \italique{maltraiter}, \italique{dévaster}. Cf., par exemple\nc{Pour d'autres exemples, qui ne manquent pas, se reporter aux \italique{testimonia} des éditions des lexiques sous-cités}, Hésychius (ε 1453 ; κ 3067) :\citComm{ἐκκοροῦσι· φθείρουσιν, ἐκκαλλύνουσιν.\\
κορέων, κορέσκων· ἐξυβρίζων.}Photius (ε 412) :\citComm{ἐκκορηθῆναι· ἐκκαλλυνθῆναι· τὸ γὰρ κορεῖν καλλύνειν καὶ σαίρειν.}Ainsi dans la \italique{Paix} (vers 58-59), Trygée s'écrie\nc{Voir Taillardat (§629) pour un exposé plus fourni sur cette image.} :\citComm{Ὦ Ζεῦ, τί ποτε βουλεύει ποεῖν;\\
Κατάθου τὸ κόρημα· μὴ ᾽κκόρει τὴν Ἑλλάδα\\
\francais{Ô Zeus, que prétends-tu faire ?\\
Dépose ton balai : ne balaie pas l'Hellade de la surface de la Terre.}}Il n'y a à notre \nieme{760} aucun sens obscène, ni là ni ailleurs\nc{Suétone (τρὶς κεκορημένε· οἱονεὶ πολλάκις ἐκσεσαρωμένος) a donc raison contre Taillardat (Περὶ βλασφημιῶν, β 260 $=$ Eustathe p.1542.47 $=$ fr. Anacréon 366 \italique{Poetae Melici Graeci} Page) et beaucoup d'autres : l'imprécation d'Anacréon `τρὶς κεκορημένε' lancée au Thrace Smerdiès signifie chastement «~triple ordure~».}. La question de Mica revient à ceci : «~Qui t'a fourragée\nc{Je m'amuse, mais comme j'appréhende déjà que cette badinerie soit prétexte à la meute libidineuse des commentamenteurs pour déchaîner leur rut philololubrique, je m'empresse de préciser : à aucun moment Critylla n'évoque l'idée que le sang qui couvre Mica puisse être hyménéen. Et d'abord parce que ce serait inepte : Mica est mère. 
Et que ceux qui tiennent tant à cette théorie exhibent seulement un passage -- ou même seulement un indice autre que des polyçonneries -- dans lequel le mot παῖδα ait pu désigner l'hymen. En attendant, il serait bon prophylactiquement d'expurger -- ἐξερᾶν, mais j'anticipe... -- du Bailly, à l'entrée ἐκκορέω, l'acception \italique{déflorer} afférente à ce passage.} ?~», «~Qui t'a ruiné [ton habit] ?~»\\
Mais ἐξεκόρησε n'est pas mis là fortuitement ; il fait entendre le mot κόρη, la «~fille~» de Mica, avec la reprise au vers suivant : τὴν ἀγαπητὴν παῖδά σου.
Le sens premier de la question de Mica est donc de savoir qui l'a ainsi salie ; le sens allusif qui l'a \italique{privée de sa fille} (ἐξ-κόρη).\\
Pour être sûr que le spectateur saisisse bien l'amphibologie, Aristophane revient à la charge au vers 761 : τίς τὴν ἀγαπητὴν παῖδά σου ᾿ξηράσατο\nc{Pour le texte, voir le commentaire critique.};
Le pronom σύ a volontairement été rejeté avant le verbe pour laisser toute ambiguïté à la crase : σοὐξηράσατο $=$ σοι ἐξηράσατο\nc{Cf. les vers 158 et 624.} d'une part ; σοὐξῃρήσατο $=$ σοῦ ἐξῃρήσατο d'autre part.\\
σοι ἐξερᾶν est le pendant de ἐκκορεῖν σε du vers précédent : Critylla demande qui a osé \italique{vider}, \italique{déverser}\nc{Peut-être même faut-il prendre ἐξερᾶν au sens de \italique{régurgiter}, \italique{recracher}, et comprendre que Mnésiloque, ajoutant l'insulte à l'outrage, après avoir empli sa bouche de vin, l'aura \italique{recraché} à la tête de Mica.} sur Mica ce que sa fille chérie avait dans le ventre.\\
σοῦ ἐξαιρεῖν : en raison de la proximité phonétique entre le `η' et le `α', le spectateur devait aisément entendre -- pour peu que l'acteur de Critylla s'y prêtât : τίς τὴν ἀγαπητὴν παῖδά σου ἐξῃρήσατο; «~Qui t'a ravi ta fille chérie ?~» ; écho au ἐκ-κόρη subrepticement ouï au vers précédent.}
